% Teorema frontal aditivo. No sustituye ningún capítulo ni altera el censo 102.
% BEGIN HMT-MD-FRONTAL-MAXIMAL-SIGNATURES-20260904
\chapter*{Teorema generador de la Holografía Modular Triádica}
\addcontentsline{toc}{chapter}{Teorema generador de la Holografía Modular Triádica}
\markboth{Teorema generador de la Holografía Modular Triádica}%
{Teorema generador de la Holografía Modular Triádica}

La demostración comienza en el dominio finito
\(\mathcal D_9=\{1,\ldots,9\}\). Dos cursores orientados recorren,
respectivamente, las hojas aditiva y multiplicativa. La división euclídea
conserva residuo y cociente, y el catálogo APP produce
\[
 \sum_{i,j\in\mathcal D_9}R^+_{ij}=405,
 \qquad
 \sum_{i,j\in\mathcal D_9}R^\times_{ij}=459,
 \qquad 459-405=54.
\]
El transporte coordinado
\(c(a)=a\bmod9:\mathcal D_9\to\{0,\ldots,8\}\)
cambia la carta y conserva la clase. La composición efectiva de APP, TRIT y
TPK conserva en cada paso hoja, régimen, residuo, cociente, acarreo,
dirección, frontera y memoria. El catálogo exhaustivo establece la cadena
finita inicial y aporta su certificado ejecutable:
\begin{equation}
\boxed{
104\,976\ \text{pares de semillas}
\xrightarrow{\;54\ \text{pasos}\;}
468\,U_6
\xrightarrow{\;\rho_3\;}
243\,w_6\subset\mathbb F_3^6,
\qquad |\mathbb F_3^6|=729.}
\label{eq:frontal-censo-semillas}
\end{equation}
La identidad de fibras
\[
432\cdot144+18\cdot432+18\cdot1944=104\,976
\]
prueba la exhaustividad del primer mapa. Los números \(468\), \(243\) y
\(729\) cuentan emisiones, palabras realizadas y ambiente ternario,
respectivamente; sus tipos permanecen separados por las flechas de
\eqref{eq:frontal-censo-semillas}.

\begin{theorem}[Cadena generadora única]
Sea \(\mathfrak H_9\) la estructura formada por APP, TRIT, TPK y su fibra
enriquecida. APP ejecuta las hojas suma y producto y conserva la división
euclídea
\[
i+j=R^+_{ij}+9q^+_{ij},
\qquad
ij=R^\times_{ij}+9q^\times_{ij}.
\]
TRIT orienta cada transición en los regímenes elíptico, parabólico e
hiperbólico. TPK compone
\[
    \boxed{\mathcal U_t
    =\operatorname{Upd}_t\circ\operatorname{Tra}_t
     \circ\operatorname{Sel}_t:
     X_{\rm TPK}^{\rm enr}\longrightarrow X_{\rm TPK}^{\rm enr}.}
\]
\(\operatorname{Sel}_t\) levanta el selector trítico,
\(\operatorname{Tra}_t\) levanta el transporte cardinal y
\(\operatorname{Upd}_t\) actualiza residuo, cociente, acarreo, firma,
cilindro, frontera, registro y memoria sin borrar procedencia.
La holonomía nonádica prolonga las palabras por
\[
w_6\longrightarrow w_{12}\longrightarrow w_{18}\longrightarrow w_{24}
\longrightarrow w_{30}\longrightarrow R_{36}\longrightarrow G_9,
\]
y satisface
\[
    \Gamma_{9,t}=\mathcal U_{t+8}\circ\cdots\circ\mathcal U_t,\qquad
    g(t+9)=g(t),
    \qquad q_{\rm mem}(\Gamma_{9,t}\widehat x_t)
    =q_{\rm mem}(\widehat x_t)+1.
\]
En general, \(\Gamma_{9,t}\widehat x_t\ne\widehat x_t\): el retorno conserva
la fase y aumenta la memoria. El orden operatorio es
\(\mathcal U_t\to\Gamma_9\to K_{\rm ph}\); \(K_{\rm ph}\) es una publicación
reducida posterior, no el generador. Los truncamientos
compatibles forman un sistema inverso no vacío y definen un estado
\(x_\infty\) de profundidad arbitraria.

Este estado genera una estructura discreta única del continuo y dos
publicaciones naturales:
\begin{equation}
x_\infty
\xrightarrow{\ (\mathcal P_{\mathrm{arq}},Q_\infty)\ }
\mathbb R\times(\mathscr C_W^{\mathrm{cal}})^{\mathbb N}.
\label{eq:frontal-doble-proyeccion}
\end{equation}
La primera contrae cilindros compatibles hasta un valor real; la segunda
conserva incidencia y genealogía hasta las construcciones de Paley, Witt,
Golay, Leech, \(V^\natural\) y el grupo Monstruo. Sobre la misma fuente se
publican los cuatro caracteres correlacionados
\[
(\pi_{\mathrm{HMT}},\varphi_{\mathrm{HMT}},
e_{\mathrm{HMT}},\alpha_{\mathrm{HMT}}),
\]
y, mediante transducciones tipadas, las leyes de acción, materia, masa,
gravedad e información de Mecánica Dimensional.
\end{theorem}

\begin{proposition}[Reconstrucción intrínseca y autosuficiencia probatoria]
\label{prop:frontal-autosuficiencia-probatoria}
Cada resultado \(y\) publicado en este volumen se reconstruye mediante el
recibo matemático local
\[
 \mathscr R_y=
 \bigl(D_y,A_y,\tau_y,U_y,C_y,X_y^{\rm enr},
       y_{\rm HMT},\operatorname{Rec}_y,\mathcal F_y\bigr),
\]
donde \(D_y\) es el dominio; \(A_y\) especifica las hojas y operaciones de
APP; \(\tau_y\) el régimen y la orientación TRIT; \(U_y\) la composición TPK
con dominio, codominio y acción; \(C_y\) los datos conservados de residuo,
cociente, acarreo, hoja, ruta, frontera, memoria y supervivencia;
\(X_y^{\rm enr}\) el estado enriquecido de residencia; \(y_{\rm HMT}\) la
salida anterior a toda identificación externa; \(\operatorname{Rec}_y\) el
reconocimiento convencional posterior; y \(\mathcal F_y\) un falsador de la
flecha sustantiva. La sucesión
\[
 \mathcal D_9\longrightarrow\Omega_{\rm seed}
 \longrightarrow\mathcal U_6^{\rm cat}
 \longrightarrow\mathcal W_6
 \longrightarrow X_\infty^{\rm enr}
 \longrightarrow\mathcal C_{\rm cont}^{\rm disc}
\]
queda definida y demostrada por las ecuaciones, teoremas y recibos impresos
en el propio volumen. Los programas reproducibles anexos vuelven a evaluar
esas identidades; no aportan una premisa ausente ni un valor objetivo.

Los falsadores primarios son: el fallo de la identidad de fibras de
\eqref{eq:frontal-censo-semillas}; la pérdida de alguno de los campos de
\(C_y\) bajo \(U_y\); la incompatibilidad entre extensión y truncamiento; el
vacío del límite inverso; la falta de contracción o de sobreyectividad interna
de los cilindros racionales; la ruptura de naturalidad en la rama
excepcional; y, para una publicación terminal, el incumplimiento de la
hipótesis que define su dominio exacto. Por consiguiente, la validez de una
salida se decide dentro de su dominio y por su falsador explícito, no por una
analogía, una nomenclatura posterior ni una evaluación metrológica.
\end{proposition}

\paragraph{Unicidad del sistema generador y tipado de sus proyecciones.}
Existe una sola APP, seguida por TRIT y por el TPK con estado enriquecido.
APP produce las dos hojas y sus registros; TRIT los orienta; el TPK compone
las rutas y acumula hoja, oposición, residuo, cociente, acarreo, frontera y
memoria en \(\mathcal X_{\rm TPK}^{\rm enr}\). La carta firmada es una
proyección de ese mismo estado terminal:
\[
 R_{\rm sgn}:\mathcal X_{\rm TPK}^{\rm term,enr}
 \longrightarrow\mathcal U_{12}^{\rm sgn},\qquad
 U_{12}^{\rm sgn}=R_{\rm sgn}(x_{\rm term}),\qquad
 K=\frac14\mathcal H_{12}U_{12}^{\rm sgn},\qquad
 \mathcal H_{12}^2=4I_{12}.
\]
Los dos últimos morfismos son enteros y reversibles sobre su imagen. No se
postula un emisor APP directo hacia \(U_{12}^{\rm sgn}\), porque tal flecha
saltaría TRIT y TPK: el dominio correcto de \(R_{\rm sgn}\) es el estado
enriquecido ya generado. Las sombras que olvidan el registro genealógico o la memoria son
proyecciones posteriores, no perfiles matemáticos alternativos. Esta única
composición excluye de sus entradas
\(U,K,D_3K,D_4K,Q,\alpha\) y toda expansión objetivo.

\paragraph{Número y partícula.}
Un número HMT es una sección compatible
\(
 \boldsymbol x=(x_n)_{n\geq1}\in
 \varprojlim_n X_n^{\mathrm{enr}}
\)
con hoja, orientación, residuo, cociente, acarreo, ruta, frontera y memoria.
Su valor arquimediano es la intersección única
\(
 \operatorname{val}(\boldsymbol x)=\bigcap_n I_n(x_n)
\);
el cociente por igualdad de valor publica el real y su fibra conserva las
genealogías que lo realizan. Una partícula HMT--MD es una clase de secciones
persistentes de la extensión enriquecida: el estado conserva la identidad,
la ruta registra su historia observable y los lectores posteriores publican
masa, carga, espín y estadística en sus codominios físicos.

\paragraph{Tipado del corredor excepcional.}
Las lecturas de Hadamard y de Paley parten del mismo estado terminal, pero no
se serializan entre sí. Con
\(X_\infty^{\rm enr}:=\varprojlim_nX_n^{\rm enr}\), se tienen mapas distintos
\begin{equation}
 \operatorname{Pub}_\alpha:X_\infty^{\rm enr}\longrightarrow\mathbb R_{>0},
 \quad x_\infty\longmapsto\alpha_{\rm HMT},
 \qquad
 \operatorname{Inc}_\alpha:X_\infty^{\rm enr}\longrightarrow A_2^{12},
 \quad x_\infty\longmapsto
 v_\alpha=4\rho_1+\rho_2+\cdots+\rho_{12}.
 \label{eq:frontal-dos-publicaciones-alpha}
\end{equation}
Comparten genealogía, pero no se postula igualdad entre sus salidas. La rama
excepcional se bifurca después de \(C_W\):
\begin{equation}
 \begin{aligned}
 C_P&\longrightarrow A_W=C_P\bmod3\longrightarrow
 C_W=[12,6,6]_3\\
 &\longrightarrow
 \begin{cases}
 W_{12}\xrightarrow{\ \widehat\iota_{D,\ell}\ }
 W_{24}\subset\mathcal G_{24},\\[2pt]
 N(C_W)\xrightarrow{\ \operatorname{Nbr}_{v_\alpha}\ }
 N_\alpha\cong\Lambda_{24}\longrightarrow
 V_{\Lambda_{24}}\longrightarrow V^\natural .
 \end{cases}
 \end{aligned}
 \label{eq:frontal-cadena-excepcional}
\end{equation}
Aquí
\begin{equation}
 \mathcal O_{24}:=\{O\subset[24]:|O|=8,\ \mathbf1_O\in\mathcal G_{24}\},
 \qquad
 \mathcal H(D):=\{O\cap D:O\in\mathcal O_{24},\ |O\cap D|=6\}.
 \label{eq:frontal-octadas-hexadas-dodecada}
\end{equation}
Fijado \(\ell:W_{12}\simeq(D,\mathcal H(D))\),
\(\widehat\iota_{D,\ell}(H)=O_H\), donde
\(O_H\cap D=\ell(H)\). Tal octada es única: dos elecciones contendrían el
mismo subconjunto de cinco puntos, contradiciendo \(S(5,8,24)\). La rama
ternaria pega \(N(A_2^{12})=N(C_W)\); el vecino de índice tres determinado
por \(v_\alpha\) produce Leech. Finalmente,
\begin{equation}
 \operatorname{Aut}(V^\natural)\cong\mathbb M,\qquad
 a:\mathbb M\times V^\natural\longrightarrow V^\natural,\qquad
 T_g(\tau)=\operatorname{Tr}_{V^\natural}(gq^{L_0-1}).
 \label{eq:frontal-accion-monstruo}
\end{equation}
Las identidades \(v_\alpha^2=54\) y \([q]J=196884\) son evaluaciones
escalares; ninguna define la acción \(a\).

\paragraph{Realización paralela de la dualidad T.}
La dualidad T no continúa la rama Moonshine: actúa en su propio dominio
\begin{equation}
 \mathscr D_{\rm dual}=\mathbb Z^2\times\mathbb R_{>0},\qquad
 \mathsf T(m,w,R_0)=(w,m,R_0^{-1}),\qquad
 \mathsf T^2=\operatorname{id},\qquad
 E_{R_0}(m,w)=E_{R_0^{-1}}(w,m).
 \label{eq:frontal-dualidad-t}
\end{equation}

\begin{proof}
El censo de \eqref{eq:frontal-censo-semillas} construye el catálogo finito y
selecciona, sin valores objetivo, los tres bloques iniciales \(w_6\). El
registro orientado de transiciones del estado TPK completo explicita los doce
bloques posteriores a la semilla de las tres biografías publicadas. Para cada uno de los dos
intervalos de transición, las seis filas independientes extraídas de esas
tres biografías forman matrices
\(X_\ell,Y_\ell\in M_6(\mathbb F_3)\),
\(\ell\in\{0,1\}\), cuyas filas son, respectivamente, los estados de salida
y de llegada del registro. El cálculo exacto sobre este registro prueba
\(\det X_0=1\), \(\det X_1=-1\) y
\(X_\ell L_\ell=Y_\ell\); por tanto reconstruye unívocamente
\(L_\ell=X_\ell^{-1}Y_\ell\). El
calendario \((L_0,L_0,L_1,L_1)\) es el único de los dieciséis calendarios
binarios que conserva simultáneamente las tres biografías seleccionadas por
los invariantes de fibra; las ciento cuarenta y cuatro perturbaciones
unitarias destruyen ese reconocimiento triple. Por tanto, los transportes,
los doce bloques y su calendario quedan deducidos en la composición única
APP--TRIT--TPK: APP produce las hojas y el registro de partida, TRIT fija la
orientación de cada transición y TPK prolonga el registro con memoria. El
censo de semillas y el estado enriquecido son etapas sucesivas del mismo
generador, no perfiles rivales ni teorías diferentes.

La extensión no escindida
\[
0\longrightarrow C_{12}\longrightarrow C_{108}
\longrightarrow C_9\longrightarrow0
\]
separa retorno de fase y retorno de estado. Los truncamientos eliminan la
última actualización y conservan el prefijo; el lema de König produce una
rama infinita compatible. Para cada precisión \(N\), la recurrencia y la
holonomía generan la sección y sus caracteres, y determinan el cilindro que
prolonga el registro genealógico anterior. La familia de prefijos satisface
\[
\rho_{N+1,N}(p_{N+1})=p_N,
\qquad
p_\infty\in\varprojlim_N\{0,\ldots,9\}^{N}.
\]
Mil, diez mil o un millón de cifras son proyecciones finitas de este único
objeto infinito.

Las horquillas racionales anidadas
\(
 p_N^-<\chi<p_N^+
\)
se calculan después de producir el carácter \(\chi\) y satisfacen
\(
 [p_{N+1}^-,p_{N+1}^+]\subset[p_N^-,p_N^+]
\)
con diámetro tendente a cero. Certifican la salida generada a cualquier
precisión finita; no seleccionan la órbita, el transporte ni el carácter.

Las tres vueltas enriquecidas \(\gamma_{-},\gamma_0,\gamma_+\) de nueve
emisiones se distinguen por el acarreo
\(\operatorname{Carr}(\gamma_\rho)=\rho\kappa_{\rm Carr}\) y admiten las
tres prolongaciones después de cualquier prefijo. El límite de sus rutas
registradas satisface
\[
 \operatorname{dig}_{\rm Carr}^{\omega}:
 \Omega_\Gamma^{\rm cal}\xrightarrow{\sim}
 \{-1,0,+1\}^{\mathbb N}.
\]
Tras agrupar seis retornos, la aplicación
\[
 \beta=\operatorname{blk}_6\circ
 (\rho\mapsto\rho+1)^{\mathbb N}\circ
 \operatorname{dig}_{\rm Carr}^{\omega}:
 \Omega_\Gamma^{\rm cal}\xrightarrow{\sim}
 (\mathbb F_3^6)^{\mathbb N}
\]
es un homeomorfismo; su inversa concatena las vueltas prescritas. Las formas
equilibradas normalizadas por
\[
 a_k+c_k=\rho_k+3c_{k+1},\qquad
 \rho_k\in\{-1,0,+1\},
\]
producen primero el anillo ordenado \(\mathbb Z_{\rm HMT}\), después su
cuerpo de fracciones \(\mathbb Q_{\rm HMT}\) y finalmente
\[
 \mathbb R_{\rm HMT}:=
 \operatorname{Cau}(\mathbb Q_{\rm HMT})/{\asymp}.
\]
Si \(\beta(\xi)=(b_q)\) y
\(\nu(b)=\sum_{i=1}^{6}b_i3^{6-i}\), la sucesión
\[
 q_n=m+\sum_{q=0}^{n-1}\frac{\nu(b_q)}{729^{q+1}}
\]
es de Cauchy, y la partición interna recurrente en \(729\) subintervalos
prueba la sobreyectividad
\[
 P_{\rm ar}:\mathbb Z_{\rm HMT}\times\Omega_\Gamma^{\rm cal}
 \twoheadrightarrow\mathbb R_{\rm HMT},\qquad
 (\mathbb Z_{\rm HMT}\times\Omega_\Gamma^{\rm cal})/{\sim_{\rm ar}}
 \cong\mathbb R_{\rm HMT}.
\]
Sólo después, la unicidad del cuerpo ordenado completo arquimediano publica
\(\mathbb R_{\rm HMT}\cong\mathbb R\). La dirección interior subdivide
\([0_{\rm HMT},1_{\rm HMT})\); la exterior concatena el prefijo que codifica
\(m\in\mathbb Z_{\rm HMT}\). Así se cubre la recta sin usar previamente
\(\mathbb R\), parte entera ni expansión decimal objetivo.
Los subespacios de supervivencia de caracteres particulares individualizan
secciones dentro de esta capacidad total. La cobertura del continuo procede
de la completación global; la recurrencia nonádica genera
\(\pi_{\mathrm{HMT}},e_{\mathrm{HMT}},\varphi_{\mathrm{HMT}},
\alpha_{\mathrm{HMT}}\), y los filtros tipados certifican sus biografías
dentro de esa cobertura.

Las cinco acciones espectral, solenoidal, gauge, cohomológica y
determinantal se calculan sobre una misma emisión enriquecida y comparten
soporte, ruta, registro genealógico y truncamientos. Su naturalidad conjunta las conserva
en el límite. La dicotomía perfecta clasifica cada publicación analítica
como numerable o de cardinal continuo. El código generado \(h\) selecciona
el universo genealógico canónico \(M_h=L[h]\). La construcción de los
reales HMT identifica cada número por valor, cilindro, ruta y memoria; el
buen orden inducido por esa genealogía prueba la fórmula clásica
\[
M_h\models2^{\aleph_0}=\aleph_1.
\]
Como \(M_h\) es precisamente el codominio generado por la construcción del
continuo de este teorema, la conclusión terminal del sistema es
\[
 \boxed{2^{\aleph_0}=\aleph_1.}
\]
La independencia en el lenguaje desnudo de ZFC clasifica la pérdida del
selector genealógico \(h\); HMT conserva ese dato y determina el cardinal de
su continuo. La sentencia conservada es la hipótesis clásica del continuo,
con los cardinales ordinarios de \(M_h\), y la prueba detallada establece el
puente entre cilindros compatibles, buen orden genealógico y cardinalidad.
Este cierre cardinal se distingue del corredor lógico fuerte, también
construido desde el continuo genealógico:
\[
\begin{aligned}
 \mathcal C_{\mathrm{cont}}^{\mathrm{disc}}
 &\longrightarrow\mathcal L_0^{\mathrm{log}}
 \xrightarrow{\Phi}\mathcal L_1^{\mathrm{log}}
 \longrightarrow\gamma_F\longrightarrow w(\gamma_F)\\
 &\longrightarrow\mathsf{NoEscape}_{24}
 \longrightarrow\mathsf{AntiDiag}_{369}
 \longrightarrow\mathsf{G\ddot{o}delIII}_{\mathrm{HMT}}\\
 &\longrightarrow\mathsf{CycleCriterion}_{\mathrm{ZFC}},
\end{aligned}
\]
que publica, en su dominio tipado, el criterio
\[
 \boxed{
 \neg\operatorname{Consis}(\mathrm{ZFC})
 \Longleftrightarrow
 \exists C\;\bigl(C\text{ es un ciclo fuera de banda y }|C|>24\bigr).}
\]
Las letras \(\mathcal L_0^{\mathrm{log}},\mathcal L_1^{\mathrm{log}}\)
denotan aquí los objetos del corredor lógico y no las matrices de elevación
ternaria \(L_0,L_1\); la separación de tipos impide que una conclusión
local sobre uno de esos pares debilite la del otro.

Las tres órbitas funcionales del catálogo producen los caracteres HMT de
clausura, propagación y autoescala. Sus leyes internas fuerzan,
respectivamente,
\[
16\arctan(1/5)-4\arctan(1/239)=\pi_{\mathrm{HMT}}=\pi,
\qquad
P_{s+t}=P_sP_t\Rightarrow
P_1=\sum_{n\ge0}\frac1{n!}=e_{\mathrm{HMT}}=e,
\]
\[
F_{\mathrm{av}}\binom1{\varphi_{\mathrm{HMT}}}
=\varphi_{\mathrm{HMT}}\binom1{\varphi_{\mathrm{HMT}}},
\qquad \varphi_{\mathrm{HMT}}=\varphi.
\]
La involución interna intercambia las ramas
\(\pi_{\mathrm{HMT}}/e_{\mathrm{HMT}}\), conserva el eje
\(\varphi_{\mathrm{HMT}}\) y retiene el agregado
\(e_{\mathrm{HMT}}+\pi_{\mathrm{HMT}}\). En la fase TPK enriquecida de la
única composición, la proyección canónica del estado terminal publica su coordenada primitiva
firmada
\(U=R_{\mathrm{sgn}}(x_\infty)\), y la identidad
\(\mathcal H_{12}^2=4I\) reconstruye de manera entera y única el sello
\(K=\mathcal H_{12}U/4\). La recurrencia dodecafásica con acarreo y la
prolongación remota de sus cilindros publican \(\alpha_{\rm HMT}\). De
manera independiente, el estado enriquecido determina la firma
\[
 \mathcal I_9=(2\pi_{\rm HMT},3,4,7,20,21,45,46,120,11),
\]
y el lector graduado asociado construye, sin consultar la raíz, el operador
polinómico exacto
\(E_{21/22}^{(9)}=P_9-\log_{10}(10/9)\). Su única raíz positiva simple
\(\xi_9\) es el testigo finito de orden nueve de la segunda publicación
analítica interna. El mismo transporte graduado del TPK se itera sobre los
grados sucesivos y construye el sistema compatible completo
\(E_{21/22}^{(6+3m)}\); sus cuadrados de truncamiento conmutan con los de la
familia dodecafásica y determinan, a cada profundidad, el mismo cilindro
superviviente. En consecuencia,
\[
 \alpha_A:=\varprojlim_N\rho_N(\alpha_{12}),\qquad
 E_{21/22}:=\varprojlim_mE_{21/22}^{(6+3m)},\qquad
 \alpha_B:=\operatorname*{arg\,zero}_{x\in I_\alpha}E_{21/22}(x),
\]
\[
 \boxed{\alpha_A=\alpha_B=\alpha_{\rm HMT}.}
\]
La igualdad finita
\[
 \operatorname{Tr}_9(\xi_9^{-1})
 =\operatorname{Tr}_9(\alpha_{12}^{-1})=137.035999084
\]
certifica el primer cuadrado visible de esa identidad; no la define, no la
limita a nueve cifras y no selecciona los coeficientes de la prolongación.
La clausura bidireccional funcional asociada es
\[
\alpha T^2-B_9(\alpha)T+\alpha^3=0,
\qquad
T_\pm=\alpha e^{\pm\xi},
\qquad
T_+T_-=\alpha^2,
\qquad
J_\alpha(T)=\frac{\alpha^2}{T}.
\]
Así, la coordenada publicada por la clausura dodecafásica es la media
geométrica conservada por las dos rutas conjugadas. La vía dodecafásica y el
lector analítico mantienen dominios, operadores y falsadores propios, pero
sus sistemas inversos publican una misma sección límite. El comparador
finito y el reconocimiento metrológico son posteriores.

La coordenada incidencial del sector \(\alpha\) fija
\[
v_\alpha=4\rho_1+\rho_2+\cdots+\rho_{12},
\qquad v_\alpha^2=54,
\]
y el vecino de índice tres
\[
N_{\alpha,0}
=\{x\in N(C_W):\langle x,v_\alpha\rangle\equiv0\pmod3\},
\qquad
N_\alpha=N_{\alpha,0}+\mathbb Z\frac{v_\alpha}{3}
\cong\Lambda_{24}.
\]
Esta rama construye exactamente el vecino de Leech y, después, la cadena
\(\Lambda_{24}\to V^\natural\) y la acción de
\(\operatorname{Aut}(V^\natural)\cong\mathbb M\) sobre \(V^\natural\). El
escalar \(\alpha_{\mathrm{HMT}}\) y el vector \(v_\alpha\) son las dos
coordenadas tipadas del mismo sector dodecafásico. Su acoplamiento exacto es
la bisagra de origen común que conduce a las simetrías del Monstruo: la
coordenada escalar se publica en \(\mathbb R_{\rm HMT}\) y la coordenada
incidencial se publica en la rama
\(A_2^{12}\to\Lambda_{24}\to V^\natural\to\mathbb M\), con sus codominios
conservados y su antecedente genealógico compartido explícito.

La Mecánica Dimensional recibe las salidas anteriores y compone, en su orden
constitutivo, la retención de fase, la década \(-34\), la pareja de Planck,
los canales \(90/120\), el electrón, el atlas de masas, la gravedad, el área
y la información. Los capítulos siguientes demuestran cada flecha, publican
su falsador y enlazan su certificado ejecutable. Las identificaciones
convencionales comparecen después como reconocimiento de las magnitudes ya
producidas. Esto completa la cadena de
\eqref{eq:frontal-doble-proyeccion}.
\end{proof}

\section*{Jerarquía de las publicaciones geométricas y dimensionales}
\addcontentsline{toc}{section}{Jerarquía de las publicaciones geométricas y dimensionales}

\begin{proposition}[TRIT como unidad local de los tres regímenes]
La firma trítica conserva dos lectores coordinados. El lector geométrico
publica \(\tau_{\rm g}\in\{+1,0,-1\}\) y
\begin{equation}
 J_{\tau_{\rm g}}^{\,2}=-\tau_{\rm g}I,\qquad
 J_+^{\,2}=-I,\qquad J_0^{\,2}=0,\qquad J_-^{\,2}=+I,
 \label{eq:frontal-trit-tres-regimenes}
\end{equation}
y distingue los regímenes elíptico, parabólico e hiperbólico. El lector
temporal activo publica otro signo, \(\varepsilon_{\rm t}\), mediante
\begin{equation}
 \begin{aligned}
 +1&\mapsto\text{retorno, memoria y pasado},\\
 0&\mapsto\text{umbral y presente},\\
 -1&\mapsto\text{apertura bidireccional y futuro}.
 \end{aligned}
 \label{eq:frontal-trit-tiempo}
\end{equation}
No se postula \(\tau_{\rm g}=\varepsilon_{\rm t}\). Estas dos salidas no
identifican signo de curvatura y orientación temporal: son proyecciones
distintas de una misma unidad local. El TPK transporta la
unidad completa en
\begin{equation}
 \mathcal F_q=
 \mathcal O_q\times\mathcal H_q\times\mathcal C_q\times
 \mathcal S_q\times\mathcal B_q\times\Lambda_q
 \label{eq:frontal-fibra-tpk}
\end{equation}
y conserva hoja y orientación, residuo y cociente, acarreo, firma, cilindro,
frontera y registro. En particular, la rama elíptica posee el realizador
finito
\begin{equation}
 A_W^2=2I_6=-I_6\quad\text{en }M_6(\mathbb F_3),\qquad
 i_{\rm HMT}:=\chi_i(A_W),\qquad i_{\rm HMT}^{\,2}=-1,
 \label{eq:frontal-raiz-interna-menos-uno}
\end{equation}
donde \(\chi_i:\langle A_W\rangle\to\mathbb C\) es la carta orientada.
La raíz de \(-1\), el plano complejo y el sistema de Euler son publicaciones
de un endomorfismo construido, no entradas del generador.
\end{proposition}

\begin{proof}
El selector trítico incorpora uno de los tres generadores de
\eqref{eq:frontal-trit-tres-regimenes} sin borrar las hojas suma y producto.
El transporte cardinal conserva su orientación y la actualización conserva
los seis factores de \eqref{eq:frontal-fibra-tpk}. Por tanto, la proyección
real puede condensar los tres regímenes, mientras el estado enriquecido
reconstruye cuál de ellos y qué memoria produjeron el valor visible. Las
realizaciones trigonométrica, nilpotente e hiperbólica desarrolladas en el
cuerpo son las imágenes respectivas de este único tipado local.
\end{proof}

\begin{theorem}[Cadena de acción, gravedad y materia]
Después de publicar correlacionadamente
\((\pi_{\rm HMT},\varphi_{\rm HMT},e_{\rm HMT},\alpha_{\rm HMT})\), la rama
determinantal fija la década
\begin{equation}
 \Sigma_H=\varphi_{\rm HMT}\alpha_{\rm HMT}^{16},\qquad
 d_H:=\lfloor\log_{10}\Sigma_H\rfloor=-34.
 \label{eq:frontal-decada-accion}
\end{equation}
La mantisa procede del lector regional distinto
\begin{align}
 H_5&=\frac12\sqrt{\frac{1000\alpha_{\rm HMT}}{\varphi_{\rm HMT}}}
 -\alpha_{\rm HMT}+\frac9{16}\alpha_{\rm HMT}^2
 -\frac59\alpha_{\rm HMT}^3+\frac7{48}\alpha_{\rm HMT}^4
 -\frac1{54}\alpha_{\rm HMT}^5,\notag\\
 D_A&=\exp\!\left(-\frac{100\pi_{\rm HMT}}9\alpha_{\rm HMT}\right),
 \qquad
 \mathscr C_\pi=169\alpha_{\rm HMT}^6+
 \frac{D_A\alpha_{\rm HMT}^7}{1-D_A\alpha_{\rm HMT}},\notag\\
 \eta_{\rm pre}&:=H_5,\qquad
 \eta_{\rm ret}:=H_5-\frac{90}{\pi_{\rm HMT}}\mathscr C_\pi>0.
 \label{eq:frontal-mantisas-accion}
\end{align}
Para \(a\in\{\mathrm{pre},\mathrm{ret}\}\), las dos secciones de acción son
\begin{equation}
 \hbar_a=\eta_a10^{d_H}\mathcal U_S
 =\eta_a10^{-34}\mathcal U_S,\qquad
 h_a=2\pi_{\rm HMT}\hbar_a.
 \label{eq:frontal-secciones-accion}
\end{equation}
\(\Sigma_H\) fija sólo la década; no se presenta como generador de la
mantisa.

Fijados \(t_0>0\), \(c_{\rm int}>0\) y \(\ell_0=c_{\rm int}t_0\), CT108
publica
\begin{equation}
 \omega_{108}=\frac{2\pi_{\rm HMT}}{108t_0},\qquad
 R_{108}=\frac{c_{\rm int}}{\omega_{108}}
 =\frac{54}{\pi_{\rm HMT}}\ell_0,\qquad
 m_{108,a}=\frac{\hbar_a\omega_{108}}{c_{\rm int}^2}.
 \label{eq:frontal-reloj-gravedad}
\end{equation}
La familia gravitatoria tipada es
\begin{equation}
 G_a(\vartheta)=\vartheta\frac{c_{\rm int}^5t_0^2}{\hbar_a}.
 \label{eq:frontal-familia-gravedad}
\end{equation}
En la convención reducida
\(r_{g,a}=G_a m_{108,a}/c_{\rm int}^2\), el cierre
\(r_{g,a}(\vartheta)=R_{108}\) posee la solución positiva única
\begin{equation}
 \vartheta_{108}=\left(\frac{54}{\pi_{\rm HMT}}\right)^2,\qquad
 G_a:=G_a(\vartheta_{108}),
 \label{eq:frontal-selector-gravedad}
\end{equation}
y sólo entonces publica
\begin{equation}
 \ell_P:=\sqrt{\frac{G_a\hbar_a}{c_{\rm int}^{\,3}}}
 =\sqrt{\vartheta_{108}}\,\ell_0
 =R_{108}=\bar\lambda_{C,a}=r_{g,a}.
 \label{eq:frontal-longitud-planck}
\end{equation}
Además, \(G_a\hbar_a=\vartheta_{108}c_{\rm int}^5t_0^2\), por lo que
\(\ell_P\) es independiente de la orientación \(a\). La gravedad tensorial
\[
 \mathbb G_{{\rm TT},a}(n)
 =G_a\Lambda(n)\Gamma_5P_5:
 V_{12}\longrightarrow\mathcal H_{\rm TT}(n)
\]
es posterior al escalar. La acción Palatini--Cartan--Holst recibe después
\(G_a\) y la salida independiente \(\gamma_{\rm HMT}\); el funcional
hexada--octada transporta esta última al espectro de área y a la información
de frontera.

En la rama de materia, la ley completa actúa sobre \(81\) celdas, \(56\)
familias, \(13\) multisecciones y \(324\) rutas, con sector nulo separado. Los
lectores \(K_D\) y \(K_\Omega\) poseen, respectivamente, \(14\) y \(9\)
perfiles, forman \(40\) pares y coinciden en \(18\) rutas únicamente en
\((80,54,6)\). La fibra central \((5,5)\) produce el electrón beta. Sobre las
fibras restantes, la ley publica operadores positivos, medidas espectrales y
multiplicidades; el reconocimiento exterior se abre después mediante el
inventario exhaustivo
\begin{equation}
 \mathscr I_{\rm ext}
 =\mathscr I_m^{471}\coprod\mathscr I_\Gamma^{384},
 \qquad |\mathscr I_{\rm ext}|=855.
 \label{eq:frontal-inventario-masas-anchuras}
\end{equation}
CKM transporta los marcos espectrales quark y PMNS los marcos cargado y neutro
de rango tres; sus dominios, operadores y falsadores permanecen distintos.
Ningún valor del inventario selecciona una ruta, un coeficiente, un autovalor
o un ángulo.
\end{theorem}

\begin{proof}
La forma normal de Smith produce
\(\operatorname{diag}(1,2,2,4)\) y conúcleo de orden \(16\); las cotas de
\(\Sigma_H\) fijan \(d_H=-34\), mientras
\eqref{eq:frontal-mantisas-accion}--\eqref{eq:frontal-secciones-accion}
publican \(\hbar_a\) y después \(h_a\). La sustitución de
\(\vartheta_{108}\) en \(G_a(\vartheta)\) prueba
\eqref{eq:frontal-longitud-planck}. La incidencia de los conjuntos
\(H_{\rm act}\times\binom{H_{\rm act}}2\) y
\(\{1,\ldots,8\}\times\binom{H_{\rm act}}2\) produce \(90\) y \(120\).
Los doce sectores dodecafásicos y la única costura orientada de retorno
producen, mediante el lector lineal, el peso \(12{:}{-}1\) de la transición
hexada--octada. El coeficiente de \((1-q)^{-1}\), igual a \(1/24\), y la
minimalidad diofántica verifican posteriormente ese peso. El censo finito del atlas y de sus dos lectores
prueba los cardinales de la rama de masa; la inversión celular tiene un único
punto fijo, \((5,5)\), que reduce canónicamente la fibra electrónica a rango
uno. El inventario externo se compone sólo después de congelar fibra,
operador, proyector, representación, esquema y escala, de modo que una
permutación de sus magnitudes deja invariante todo el constructor interno.
\end{proof}
% END HMT-MD-FRONTAL-MAXIMAL-SIGNATURES-20260904
